\documentclass[11pt]{article}
\usepackage[margin=2.6cm]{geometry}
\usepackage{mathpazo}
\usepackage{microtype}
\usepackage{booktabs}
\usepackage{graphicx}
\usepackage{float}
\usepackage{xcolor}
\usepackage{titlesec}
\usepackage{caption}
\usepackage{listings}
\usepackage[hidelinks]{hyperref}
\emergencystretch=3em

\definecolor{accent}{HTML}{1F4E79}
\definecolor{good}{HTML}{1A7F37}
\definecolor{bad}{HTML}{B42318}
\titleformat{\section}{\Large\bfseries\color{accent}}{\thesection}{0.8em}{}
\titleformat{\subsection}{\large\bfseries\color{accent!80!black}}{\thesubsection}{0.8em}{}
\captionsetup{font=small,labelfont={bf,color=accent}}
\newcommand{\armA}{\textcolor{bad}{arm A}}
\newcommand{\armB}{\textcolor{accent}{arm B}}

\lstdefinestyle{json}{
  basicstyle=\ttfamily\footnotesize,
  backgroundcolor=\color{gray!8},
  frame=single, rulecolor=\color{gray!40}, framerule=0.4pt,
  breaklines=true, columns=fullflexible, keepspaces=true,
  showstringspaces=false, aboveskip=8pt, belowskip=8pt}
\lstset{style=json}

\newsavebox{\eliBox}
\newenvironment{eli5}
  {\par\medskip\noindent\begin{lrbox}{\eliBox}\begin{minipage}{\dimexpr\linewidth-2\fboxsep-2\fboxrule}\small}
  {\end{minipage}\end{lrbox}\noindent\fcolorbox{accent!40}{accent!4}{\usebox{\eliBox}}\par\medskip}

\title{\color{accent}\bfseries Structured DAG-Diff Mutation for CARL Reasoning Chains\\
\large The Diff Language, Its Grammar Enforcement, and the A/B Evaluation\\
\large (aligned100 --- evidence-field schema parity with the standard operator)}
\author{GigaEvo}
\date{2026-07-03}

\begin{document}
\maketitle

\begin{abstract}
\noindent
We replaced free-form JSON rewriting of CARL reasoning-chain genomes with a
\emph{schema-constrained positional-slot diff}: the mutation LLM emits a typed edit whose
grammar is compiled per parent set, validated with pydantic, and applied deterministically.
This run adds one change on top of the shipped grammar: the diff schema is now
\emph{aligned} with the standard mutation operator --- it emits the same evidence surface
(\texttt{archetype}, \texttt{justification}, \texttt{insights\_used},
\texttt{insight\_ids\_used}, \texttt{card\_ids\_used}, \texttt{changes}), replacing the old
single free \texttt{reasoning} field, so the two operators cite evidence and hypothesise
change identically (\S\ref{sec:alignment}). In a paired 100-mutant A/B on the summarizer
task, the diff arm completed with \textbf{zero json-parse failures and zero structurally
invalid programs persisted among completed evaluations}; the rewrite arm lost 3 mutants
(all \texttt{json\_parse\_error}) to malformed children that each burned a full evaluation.
Both arms took the same 4 transient pre-persist LLM failures. Fitness is at parity (diff
child-mean marginally higher, rewrite peak higher --- a single replicate, no variance floor).
The alignment moved the token bill: the evidence fields raised the diff arm's output tokens
$+16\%$ over the pre-alignment diff, so output is now \emph{above} the rewrite arm
($+8.5\%$), while input stays $+61\%$ from shipping the schema every call --- the diff is no
longer output-cheaper, it buys evidence parity and zero parse failures at a token premium.
The instrumentation fires on 100/100 diff calls; under \texttt{memory=none} it reads
\texttt{0/0} insights grounded (nothing to cite) --- it will carry signal only with memory
on. Two earlier diagnostic runs exposed the last two grammar gaps (a validator-checked
ordering rule, then unbounded dependency arrays); both closures are documented as design
history (\S\ref{sec:encodings}, \S\ref{sec:depfix}).
\end{abstract}

\section{Motivation}
CARL chain genomes are wire-JSON documents (\texttt{ReasoningChain.to\_dict} plus platform
extras). The stock mutation path asks the LLM to rewrite the whole document as free text:
every emission re-risks the entire syntax, and a malformed child is only discovered
\emph{after} it has been persisted and has consumed a full evaluation. The hypothesis:
constraining mutation to a typed diff language --- where structural validity is enforced by
the provider's constrained decoding and a pydantic contract --- eliminates that failure class
at the cheapest possible stage.

\section{The diff language}
\label{sec:language}
This section describes the language in its shipped form --- the fixed-key \emph{object
encoding} --- which is exactly what the A/B in \S\ref{sec:experiment} ran. How the design
arrived here (two earlier diagnostic runs, each closing one grammar gap) is told in
\S\ref{sec:encodings} and \S\ref{sec:depfix}.

\subsection{ELI5}
\begin{eli5}
Imagine the child chain as a shelf of eight numbered boxes. To describe a new chain, you
fill boxes from the left. Into each box you put either a \emph{photocopy} of a step from any
parent --- naming it by its printed label, like \texttt{a2}, optionally with corrections
scribbled on top --- or a \emph{brand-new} step you write yourself. Every box also declares
which earlier boxes it reads its input from, and here is the trick: the form for box $k$
only has checkboxes for boxes $1..k-1$. Pointing at yourself or at a later box is not
``against the rules'' --- there is literally no place on the form to write it. Box~1 has no
checkboxes at all; it always reads the raw task input. When you are done, the remaining
boxes stay empty (no gaps allowed), and whatever sits in the last filled box produces the
answer that gets graded.
\end{eli5}
\noindent
The LLM never writes a chain; it fills in this form. The form is a JSON schema compiled
fresh for each mutation from the actual parents, and the provider's constrained decoding
guarantees the reply matches the form --- a violating token can never even be sampled.

\subsection{Anatomy of a diff}
A diff \emph{is} the full child chain. One payload, real shape (parents rendered as
\texttt{a1..a2} and \texttt{b1..b3}):
\begin{lstlisting}
{
  "archetype": "Guided Innovation",
  "justification": "Insight (faithfulness, final step): the polish
      step paraphrases away source wording the metric scores. Insert
      a fact-check between draft and polish so the draft is checked
      against the source before polishing.",
  "insights_used": ["final step drops source wording under paraphrase"],
  "insight_ids_used": ["a::struct::0"],
  "card_ids_used": [],
  "changes": ["insert a new Verify-facts step between a1 and a2",
              "rewire a2 to read the verified draft"],
  "base_parent": "A",
  "slot_1": {"kind": "keep", "id": "a1"},
  "slot_2": {"kind": "new",
             "title": "Verify facts",
             "aim": "Cross-check the draft against the source text;
                     list any unsupported claim.",
             "dependencies": ["slot_1"]},
  "slot_3": {"kind": "keep", "id": "a2",
             "edits": {"aim": "Rewrite the draft into one crisp
                       sentence, keeping verified source wording."},
             "dependencies": ["slot_2"]},
  "slot_4": null, "slot_5": null, "slot_6": null,
  "slot_7": null, "slot_8": null
}
\end{lstlisting}
Field by field:
\begin{itemize}
  \item \textbf{evidence surface} (\texttt{archetype}, \texttt{justification},
  \texttt{insights\_used}, \texttt{insight\_ids\_used}, \texttt{card\_ids\_used},
  \texttt{changes}) --- identical to the standard mutation operator's output minus its
  \texttt{code} field, whose role the positional slots play (\S\ref{sec:alignment}).
  \texttt{archetype} names the move from the shared vocabulary; \texttt{justification} is a
  falsifiable causal hypothesis (prompt rules forbid tautologies and diff restatements);
  \texttt{insight\_ids\_used}/\texttt{card\_ids\_used} ground the edit in the parent's
  structured insight ids and injected memory-card ids, which the citation-integrity check
  scores; \texttt{changes} enumerates the concrete edits.
  \item \textbf{\texttt{base\_parent}} --- an enum of the parent letters actually present.
  This is \emph{lineage attribution}: the parent whose run metadata and task description the
  child inherits, and the lineage credit-assignment target. It does \emph{not} restrict which
  steps may be kept.
  \item \textbf{\texttt{slot\_1} \dots\ \texttt{slot\_8}} --- the child's steps, in execution
  order. \texttt{slot\_1} is required; the rest admit \texttt{null}. Filled slots must form a
  prefix (no gaps). The last filled slot is the output step --- its answer is what gets
  scored.
\end{itemize}

\subsection{The two slot forms}
Each filled slot is one of two typed objects, discriminated by \texttt{kind}:
\begin{description}
  \item[\texttt{keep}] Reuse a rendered parent step by id. \texttt{id} is an enum compiled
  from the actual parents --- parent-prefixed (\texttt{a1}, \texttt{b3}, \dots), spanning
  \emph{all} parents, so a keep may pull a donor's step directly (crossover is free).
  Optional \texttt{edits} overlays new values onto any content field
  (\texttt{title}, \texttt{aim}, \texttt{stage\_action}, \texttt{reasoning\_questions});
  unset fields keep the parent's text. \texttt{title}/\texttt{aim} edits must be non-empty
  --- the schema carries the executor's own invariant.
  \item[\texttt{new}] A fresh LLM step: required non-empty \texttt{title} and \texttt{aim};
  optional \texttt{stage\_action} and \texttt{reasoning\_questions}.
\end{description}
Both forms take the same \texttt{dependencies} field --- except in \texttt{slot\_1}, where
the field does not exist.

\subsection{Wiring: the position-narrowed dependency enum}
\texttt{dependencies} lists the earlier slots whose outputs this step consumes --- slot
positions in the \emph{new} chain, never parent ids. The schema for slot $k$ offers the enum
$\{\texttt{slot\_1},\dots,\texttt{slot\_{k-1}}\}$ and nothing else:
\begin{center}
\footnotesize
\begin{tabular}{ll}
\toprule
slot & \texttt{dependencies} may mention \\
\midrule
\texttt{slot\_1} & --- (field absent; reads the raw task input) \\
\texttt{slot\_2} & \texttt{slot\_1} \\
\texttt{slot\_3} & \texttt{slot\_1}, \texttt{slot\_2} \\
$\vdots$ & $\vdots$ \\
\texttt{slot\_8} & \texttt{slot\_1} \dots\ \texttt{slot\_7} \\
\bottomrule
\end{tabular}
\end{center}
An empty list is legal anywhere it exists --- a step may read the raw input in parallel with
slot 1. The array is also \emph{length-capped} in the grammar (\texttt{maxItems}: $k-1$ for
slot $k$): a step can never cite more distinct predecessors than exist, and a sampler that
starts looping on a dependency token hits a forced \texttt{]} after at most seven items
instead of running away. Acyclicity of the child DAG is therefore a \emph{typing} fact, not
a checked property: every edge points strictly backwards by construction.

\subsection{Every mutation move, spelled in the language}
\begin{center}
\small
\begin{tabular}{ll}
\toprule
move & spelling \\
\midrule
edit a step & \texttt{keep} it with \texttt{edits} overriding fields \\
insert a step & put a \texttt{new} slot between two \texttt{keep}s, rewire \\
delete a step & simply never fill a slot with it \\
reorder & fill the slots in the new order \\
duplicate & \texttt{keep} the same id in two slots \\
rewire only & same \texttt{keep}s, different \texttt{dependencies} \\
crossover & \texttt{keep} ids from more than one parent \\
full rewrite & every slot \texttt{new} \\
branching / parallel DAG & two slots share a dependency; a later slot joins them \\
\bottomrule
\end{tabular}
\end{center}
Two of these in miniature. Deletion is pure omission --- collapsing a 2-step parent to its
final step is a one-slot diff:
\begin{lstlisting}
{"archetype": "Harmful Pattern Removal", "justification": "...",
 "changes": ["drop the draft step, keep only the polish"],
 "base_parent": "A",
 "slot_1": {"kind": "keep", "id": "a2"},
 "slot_2": null, ..., "slot_8": null}
\end{lstlisting}
Crossover is just a keep-id from another parent --- graft B's fact-extractor in front of A's
polisher:
\begin{lstlisting}
{"archetype": "Approach Synthesis",
 "justification": "... parent:B step:b1 ...",
 "changes": ["prepend b1 fact-extractor before a2 polish"],
 "base_parent": "A",
 "slot_1": {"kind": "keep", "id": "b1"},
 "slot_2": {"kind": "keep", "id": "a2", "dependencies": ["slot_1"]},
 "slot_3": null, ..., "slot_8": null}
\end{lstlisting}

\subsection{What the grammar makes impossible}
The design goal: a defect class should be \emph{unrepresentable}, not detected. What an
unconstrained rewriter routinely gets wrong, this language cannot even spell:
\begin{table}[H]\centering\small
\begin{tabular}{ll}
\toprule
illegal child & why it cannot be emitted \\
\midrule
self-dependency (\texttt{slot\_k} $\to$ \texttt{slot\_k}) & slot $k$'s enum stops at \texttt{slot\_{k-1}} \\
forward dependency (\texttt{slot\_2} $\to$ \texttt{slot\_5}) & same enum narrowing \\
any dependency on slot 1 & the field does not exist there; extra keys rejected \\
dependency cycle & impossible transitively: all edges point backwards \\
keep of a nonexistent step (\texttt{a9}) & \texttt{id} enum is compiled from the actual parents \\
9-step chain & there is no \texttt{slot\_9} property \\
empty chain & \texttt{slot\_1} is required \\
empty \texttt{title}/\texttt{aim} (new or edit) & \texttt{minLength: 1} in the schema \\
runaway \texttt{dependencies} list (\texttt{slot\_1} spammed 70$\times$) & \texttt{maxItems}: $k-1$ caps the array \\
invented fields, malformed JSON & \texttt{additionalProperties: false}; constrained decoding \\
\midrule
gap fill (\texttt{slot\_3} filled, \texttt{slot\_2} null) & \emph{the} residual pydantic validator (contiguity) \\
\bottomrule
\end{tabular}
\caption{Everything above the line is enforced by the decoding grammar itself. Contiguity is
the single rule left to a validator --- and combined with the enums it also makes
``dependency on an empty slot'' impossible.}
\end{table}

\subsection{How enforcement works}
\begin{enumerate}
  \item \textbf{Per-mutation compilation.} \texttt{AllowedDagChanges.build\_schema} parses
  the actual parents and generates a pydantic model with \texttt{create\_model}: the keep-id
  enum, the \texttt{base\_parent} enum, and eight per-position slot unions are all derived
  from what really exists right now.
  \item \textbf{Constrained decoding.} The JSON schema ships to the provider as the
  structured-output contract (\texttt{method="json\_schema"}). The provider compiles it into
  a decoding grammar; tokens that would violate the schema get zero probability. The model
  cannot emit a forward reference for the same reason it cannot emit an unbalanced brace.
  \item \textbf{Portable subset.} Provider grammar compilers accept different JSON-schema
  dialects; Gemini is the strictest we probed (hard-rejects \texttt{\$ref}/\texttt{\$defs}
  and \texttt{const}). \path{gigaevo/llm/schema_compat.py} emits every schema in the
  lowest-common-denominator subset: refs inlined, \texttt{const} $\to$ single-value
  \texttt{enum}, decorative keys dropped. Nothing is provider-conditional at runtime.
  \item \textbf{Grammar budget.} Strict compilers limit construct complexity, not bytes (a
  padded 120KB schema passes; an 87KB slot-union grammar does not). Branching the slot
  models per parent blew the cliff at 2 parents $\times$ 8 slots ($\sim$27KB); the shipped
  encoding uses a \emph{single} model --- one global parent-namespaced keep-id enum,
  \texttt{base\_parent} as a plain enum field, no union discriminator --- at 13.7KB with
  $\sim$2$\times$ headroom.
  \item \textbf{Pydantic backstop.} The same model validates the reply before anything
  touches the engine; the contiguity rule lives here. A rejection costs one LLM call ---
  nothing is persisted, nothing is evaluated.
\end{enumerate}

\subsection{From diff to child}
The applier is deterministic --- the LLM's creative surface ends at the diff:
transcription walks \texttt{slot\_1..slot\_8} until the first \texttt{null}; \texttt{keep}
slots copy content fields from the referenced parent step (any parent) and overlay
\texttt{edits}; \texttt{dependencies} become sorted, deduplicated integer positions; run
metadata and the task description come from \texttt{base\_parent}; the last step is stamped
as the output step. The result is round-tripped through \texttt{ReasoningChain} as a final
assertion and emitted as canonical wire JSON --- the child is indistinguishable in format
from a hand-written genome.

\section{Implementation}

\subsection{Operator and vocabulary seam}
Two pieces, both behind existing seams:
\begin{itemize}
  \item \path{gigaevo/evolution/mutation/structured_diff.py} --- the generic engine-level
  operator \texttt{StructuredDiffMutationOperator}: it asks an \texttt{AllowedChanges}
  vocabulary to build a \texttt{DiffSchema} for the current parents, runs a
  structured-output LLM call against that schema, validates, and applies. It knows nothing
  about chains.
  \item \path{gigaevo/chains/dag_changes.py} --- \texttt{AllowedDagChanges}, the
  chain-specific vocabulary of \S\ref{sec:language}: schema compilation, parent rendering
  (the \texttt{a1}/\texttt{b2} listings the LLM sees), the applier, and the natural-language
  contract injected into the mutation prompt.
\end{itemize}

\subsection{The encoding search}
\label{sec:encodings}
Getting the diff grammar through the strictest compiler took four encodings:
\begin{table}[H]\centering\footnotesize
\begin{tabular}{llll}
\toprule
encoding & ordering rule in grammar? & size & Gemini verdict \\
\midrule
E1: tuple-union (\texttt{prefixItems}) & yes (per-slot models) & $\sim$87KB & rejected (complexity) \\
E2: cons-list recursion (\texttt{\$ref}) & yes (per-depth) & --- & rejected (\texttt{\$ref} banned) \\
E3: flat \texttt{anyOf[keep,new]} array & \textcolor{bad}{no --- pydantic validator} & $\sim$8KB & accepted --- superseded \\
E4: fixed-key object (\texttt{slot\_k}) & \textcolor{good}{yes (narrowed enums + \texttt{maxItems})} & 13.7KB & \textcolor{good}{accepted --- \textbf{shipped, ran the A/B}} \\
\bottomrule
\end{tabular}
\caption{Raw size is not the limit; construct complexity is. E4 is the encoding described in
\S\ref{sec:language} and the one the A/B ran; E3 differed only in expressing slots as an
array and checking the ordering rule post-decode (\S\ref{sec:depfix}).}
\end{table}

\subsection{Engine integration}
The engine assumed the mutator's \texttt{base\_parent} is a 1-based integer (wire-JSON
convention); diff payloads use namespace letters, so every diff mutant died pre-persist in
the first smoke (\texttt{int('A')}). Fixed test-first with \texttt{base\_parent\_index()} in
\texttt{gigaevo/evolution/engine/mutation.py} (int or namespace letter; garbage falls back to
parent 1 with a warning). Re-reviewing the aligned schema surfaced a \emph{second} consumer
that had not routed through this helper: \texttt{IntraMemoryStage}'s stamp-less fallback for
a crossover child's \texttt{crossover\_role} label read \texttt{base\_parent} with an
\texttt{isinstance(int)} gate, silently mis-defaulting the letter \texttt{"B"} to parent~0
and flipping the base/donor attribution the analyst prompt depends on. Pre-existing and off
the normal path (the base-id stamp populates upstream), but every diff-operator crossover
child is exactly the letter case --- fixed test-first to route through
\texttt{base\_parent\_index} as well. Other pre-launch fixes (config-group promotion, loader
pattern clobber, prompt-fairness rewrite for arm A, token accounting) are in
\texttt{04\_issues\_log.md}.

\section{Evidence-field alignment}
\label{sec:alignment}
Before this run the diff operator emitted a single free \texttt{reasoning} string, while the
standard \texttt{MutationAgent} emits a structured evidence surface:
\texttt{archetype} (from a shared vocabulary), \texttt{justification},
\texttt{insights\_used}, \texttt{insight\_ids\_used}, \texttt{card\_ids\_used}, and
\texttt{changes}. The two operators hypothesised change in different shapes, so their
children carried non-comparable self-reports and the diff children were invisible to the
citation-integrity check.

The alignment makes the diff schema mirror \texttt{MutationStructuredOutput} \emph{minus its
\texttt{code} field} --- whose role the positional slots already play. Field descriptions are
read from the standard model via a shared mirror so the two schemas cannot drift; both
structured\_diff prompts (generic and the problem-local summarizer) were rewritten
section-for-section to match their mutation counterparts, differing only in the output-form
section (diff language vs.\ Python / CARL wire JSON) and the letter-vs-number parent
citation. Citation integrity and the archetype/justification metadata now flow through the
diff agent and operator on the same path as the standard operator.

\paragraph{Run-level verification.} Every one of \armB's 100 persisted children carries the
new \texttt{mutation\_output} surface: \texttt{justification} on 100/100, a \texttt{changes}
list on 51/100, and \texttt{archetype} drawn entirely from the shared vocabulary (Approach
Synthesis 39, Precision Optimization 35, Guided Innovation 19, Solution Space Exploration 3,
Harmful Pattern Removal 3, Proven Pattern Extension 1). Insight grounding fired: 51 children
cited 2--7 parent-insight ids, the rest zero. The citation-integrity check ran on 100/100
diff calls --- but under \texttt{memory=none} there are no injected insights or cards to
cite, so it reads \texttt{0/0 grounded} throughout: the instrumentation is live and correct,
and will carry signal only in a memory-on run. \texttt{card\_ids\_used} is 0 everywhere,
consistent with \texttt{0/0}. The sampled child's \texttt{base\_parent} is the letter
\texttt{"B"} --- exactly the encoding the \texttt{IntraMemoryStage} fix above now resolves.

\section{Experiment}
\label{sec:experiment}
Paired 100-mutant runs, identical seeds (4), engine config, and models: mutation
gemini-3.5-flash (OpenRouter, reasoning effort high); chain executor Qwen3-235B-Instruct
(LiteLLM proxy). Task: \texttt{problems/chains/summarizer}, ROUGE-L fitness,
\texttt{memory=none}.
\textbf{\armA{}} = stock full wire-JSON rewrite with a chain-faithful prompt;
\textbf{\armB{}} = structured slot diff on the shipped object encoding E4 of
\S\ref{sec:language}, \texttt{maxItems} included, now emitting the \emph{aligned} evidence
surface of \S\ref{sec:alignment}.

\section{Results}

\subsection{Validity funnel (headline)}
\begin{table}[H]\centering
\begin{tabular}{lcc}
\toprule
 & \armA{} (rewrite) & \armB{} (diff) \\
\midrule
mutation LLM calls & 107 & 107 \\
generation failures (pre-persist, transient LLM) & 4 (no code extracted) & 4 (llm call error) \\
children persisted & 100 & 100 \\
in flight at the 100-mutant stopper (excluded) & 6 & 8 \\
children with a completed evaluation & 94 & 92 \\
\textbf{\texttt{json\_parse\_error} (invalid, completed)} & \textbf{\textcolor{bad}{3}} & \textbf{\textcolor{good}{0}} \\
diff schema / apply rejections & --- & \textbf{\textcolor{good}{0}} \\
valid children (completed) & 91 & 92 \\
valid yield per completed evaluation & 96.8\% & \textbf{100\%} \\
\bottomrule
\end{tabular}
\caption{The arms fail at different stages. The ``in flight at the stopper'' row is an
artifact of the fixed mutant budget, not a failure: when the 100th child persists, in-flight
evaluations are swept; those children parse as valid chains and are excluded from the yield
denominators. Both arms took the same 4 transient pre-persist LLM failures (OpenRouter/proxy
contention).}
\end{table}
\noindent
\armA{} defects pass mutation silently (the operator extracts a string) and die in
evaluation: 3 mutants persisted malformed JSON (truncated or badly quoted), each burning a
mutant slot plus a full evaluation pipeline before being marked invalid. \armB{} recorded
\textbf{zero} \texttt{json\_parse\_error} and \textbf{zero} schema/apply rejections in 107
calls: with the wiring rules in the grammar and the array caps of \S\ref{sec:depfix},
nothing structurally invalid can reach the population, by construction. The prior run's
degenerate self-reference / repetition rejections did not recur.

\subsection{Fitness --- parity}
Child mean 0.392 (\armA) vs.\ 0.413 (\armB); medians 0.398 vs.\ 0.419 --- the diff arm's
child distribution sits marginally higher (a tighter, more reliable edit distribution is the
expected signature of small schema-constrained moves). But the rewrite arm owns the peak:
its best child is 0.592 against \armB's 0.497 (whose overall best is essentially its own seed
at 0.495). The same 4 seeds re-score 0.489 (\armA) vs.\ 0.495 (\armB) across the two runs
purely from executor sampling noise, and no same-config variance floor exists for this
problem yet; with a single replicate per arm, \armA's lone 0.592 child is a one-off top
sample. We claim parity --- a marginal mean edge to the diff, a peak to the rewrite, neither
separable from noise.
\begin{figure}[H]\centering
\includegraphics[width=\textwidth]{ab_overview.png}
\caption{Left: per-child fitness (dots) and best-so-far (line). Right: validity funnel.}
\end{figure}

\subsection{Structural exploration}
Both arms span 1--3 steps in this run: \armA{} $\{29, 44, 24\}$ across lengths 1/2/3,
\armB{} $\{18, 61, 20\}$ --- \armB{} concentrates harder on 2-step chains. The language
admits up to 8 slots and an earlier diagnostic run did produce 4- and 8-step children, but
the effect did not replicate here, so we claim no structural-exploration advantage; the
grammar's guarantee is that whatever lengths arrive are well-formed.
\begin{figure}[H]\centering
\includegraphics[width=0.7\textwidth]{ab_nsteps.png}
\caption{Chain-length distribution of valid children.}
\end{figure}

\subsection{Cost}
\begin{table}[H]\centering
\begin{tabular}{lccc}
\toprule
mutation side (gemini-3.5-flash) & \armA{} & \armB{} & \armB{} vs \armA{} \\
\midrule
tokens in & 755{,}173 & 1{,}213{,}975 & $+61\%$ \\
tokens out & 766{,}831 & 832{,}204 & \textbf{$+8.5\%$} \\
mean output tokens per completed call & $\sim$7{,}670 & $\sim$8{,}490 & $+11\%$ \\
\bottomrule
\end{tabular}
\caption{The alignment moved the bill. Pre-alignment (terse \texttt{reasoning}) the diff arm
was output-\emph{cheaper} than the rewrite; adding the evidence surface raised \armB{}'s
output $\sim$16\% (714{,}864 $\to$ 832{,}204), so output is now \emph{above} \armA{}
($+8.5\%$). Input stays a large premium ($+61\%$), and the driver is the \emph{schema}, not
the prompt: both arms render both parents (\texttt{num\_parents}$=$2) and \armB's parent
\emph{digest} is in fact smaller than raw wire JSON, so the parent content is common-mode.
What differs is the constrained-decoding contract itself --- the diff schema is
$\sim$4.4$\times$ the standard operator's (16.9\,KB vs.\ 3.8\,KB, its per-position enums
narrowed by both parents' step-ids), charged as input on every call. That $\sim$13\,KB
schema differential accounts for essentially the whole $\sim$4.4k-token input gap; the
prompt text is comparable (system prompts within 0.7\,KB). Output tokens include the model's
reasoning channel (effort high), which
varies run to run. Executor completion tokens: 80K (\armA) vs.\ 86K (\armB). Net: the diff
no longer buys a token saving --- it buys evidence parity and zero parse failures at a token
premium.}
\end{table}
\begin{figure}[H]\centering
\includegraphics[width=\textwidth]{ab_tokens.png}
\caption{Per-mutant token spend. Left: output tokens per mutation call (dots) with a 10-call
rolling mean (line). Zero-token points are calls that returned no completion (\armA's
extraction failures, or in-flight calls swept when the mutant budget was reached). Right:
cumulative mutation tokens; the two output curves now run close together (the diff's former
output saving is gone once the evidence fields are emitted), while \armB's input curve sits
well above \armA's from shipping the schema with every call.}
\end{figure}

\section{Design history: two diagnostic runs, two grammar gaps closed}
\label{sec:depfix}
The grammar the A/B ran was hardened by two earlier 100-mutant diagnostic runs, each of
which exposed exactly one failure class and forced it into the schema. The requirement was
explicit throughout: fix defects in the grammar, not post-hoc.

\subsection{Gap 1: the ordering rule lived in a validator (encoding E3)}
The first run used E3 --- slots as a flat array, the \emph{dependencies-point-backwards}
rule checked by pydantic after decoding. All three of its pre-persist rejections (3/108
calls) were the same degenerate violation: slot~1 emitting
\texttt{dependencies:\ ["slot\_1"]}, a self-reference on the first slot --- precisely the
one rule the grammar did not carry. The fix is the fixed-key object encoding of
\S\ref{sec:language}: \texttt{slot\_1..slot\_8} as separate properties, each slot's
dependency enum narrowed to earlier slots, slot~1 with no dependencies field at all.
Self- and forward-references became unrepresentable; the sole surviving validator is
contiguous fill. Probes shaped the final form: Gemini rejects the per-parent-branched
variant (grammar cliff between 7 and 8 slots) and nested union discriminators, but accepts
the single-model variant --- keep-ids in one global parent-namespaced enum,
\texttt{base\_parent} a plain enum --- at 13.7KB with $\sim$2$\times$ headroom. A welcome
side effect: keeps from \emph{any} parent are now first-class, so crossover no longer
requires re-creating donor steps. Design notes:
\texttt{experiments/carl\_dag\_diff\_dependency\_rule\_fix.md}.

\subsection{Gap 2: unbounded arrays outrun the provider's enum enforcement}
The second run, on the object encoding, produced zero ordering violations --- and two
rejections of a class nobody had predicted: \emph{degenerate repetition}. The sampler
locked into a loop and emitted \texttt{"slot\_1"} $\sim$70 times inside one
\texttt{dependencies} array; deep into the list, items came back as \texttt{"slot"} and
\texttt{""} --- the provider's enum enforcement demonstrably lapses far into a long array.
The pydantic backstop caught both replies (nothing was persisted), but the calls were
wasted. The grammar fix: \texttt{maxItems}: $k-1$ on slot $k$'s dependency array
(\texttt{Field(max\_length=k-1)}). A repetition loop now hits a forced \texttt{]} after at
most seven items --- well inside the enforcement horizon --- and a step can never cite more
distinct predecessors than exist. The A/B of \S\ref{sec:experiment} ran with this cap and
recorded \textbf{zero} schema errors in 107 calls.

\subsection{Verification}
\begin{itemize}
  \item \textbf{Unit:} 32 tests in \path{tests/chains/test_dag_changes.py} --- every mutation
  archetype constructs a valid child; eleven illegal-payload classes (self-dep, forward dep,
  slot-1 dep, dangling id, gap fill, empty chain, empty aim, 9th slot, over-cap dependency
  list, the 71-item degenerate repetition payload itself, \dots) are rejected; schema-shape
  tests walk the emitted JSON schema and assert the per-position enum narrowing and the
  \texttt{maxItems} caps; a 300-case fuzz applies random schema-valid diffs and round-trips
  every child.
  \item \textbf{Suite:} the aligned-schema scope --- \path{tests/chains}, the diff agent,
  prompt$\leftrightarrow$schema reconciliation, citation integrity, the diff operator, the new
  \path{tests/evolution/test_diff_output_parity.py} field-parity guard, and the ideas-tracker
  models --- passed 400/400 (2 skipped); the \texttt{base\_parent} letter-recovery fix in
  \path{lineage_memory.py} carries its own regression test; lint clean.
  \item \textbf{Live smoke} (5 mutants, gemini-3.5-flash, same launch path as the A/B):
  5/5 mutation calls grammar-accepted, zero schema rejections, all children valid ---
  confirming Gemini accepts \texttt{maxItems} before the full run.
\end{itemize}

\section{Verdict}
The hypothesis is \textbf{supported end-to-end on the shipped grammar, now with evidence
parity}: \texttt{json\_parse\_error} invalid programs $3\to0$, valid yield per completed
evaluation $96.8\%\to100\%$, zero schema/apply rejections in 107 calls, fitness at parity
(a marginal child-mean edge to the diff, the peak to the rewrite, neither separable from
noise). The evidence-field alignment lands: \armB's children now carry the same
\texttt{archetype}/\texttt{justification}/\texttt{insights\_used}/\texttt{changes} surface as
the standard operator, on 100/100 children, and the citation-integrity check runs on every
diff call (trivially \texttt{0/0} only because \texttt{memory=none} injects nothing to
cite). The cost picture changed with it: the diff is no longer output-cheaper --- the
evidence fields push output to $+8.5\%$ over the rewrite, against a standing $+61\%$ input
premium from shipping the schema every call. Every wiring rule of the diff language remains
unrepresentable-to-break, and the only validator left is contiguity. Next: a memory-on run
to give the citation-integrity signal something to measure, and longer runs to test whether
the 8-slot headroom converts to structural exploration and fitness.
\end{document}
